% Example content adapted from David Molik’s tenure document.
\PortfolioSection{Summary of Candidate’s Non-Directed Service}
\SectionGuidance{Describe service beyond your assigned responsibilities to the Libraries, university, profession, and public. Explain your role and the contributions or outcomes of each activity.}
\ExampleNote

\subsection*{Service to Kansas State University Libraries}

During the evaluation period, I contributed non-directed service to Kansas State University Libraries through committee participation, professional development, and representation in national forums that directly support library priorities and capacity.

I participated in advanced leadership and assessment training through Harvard’s Leadership Institute for Academic Librarians (LIAL), the Kansas State University Fall New Department Head Institute, and the University of Kansas Assessment in Action Workshop. These activities strengthened my capacity in leadership, assessment, and evidence-based decision-making and informed my service and leadership within the Libraries.

I served as a member of the Environmental Health and Safety (EHS) Committee, contributing to institutional oversight related to research and workplace safety, and as a member of the Wildcat Scholar Committee, where I provided scholarly communication and research infrastructure expertise to discussions related to research visibility, repository integration, and institutional research profiling.

\subsection*{Service to Kansas State University}

Through my committee service and liaison activities, I contributed to broader University initiatives related to research infrastructure, compliance, and governance. My participation on the Wildcat Scholar Committee supported institution-wide efforts to improve the integration of research outputs, profiles, and repository services, strengthening Kansas State University’s research visibility and reporting capacity.

In addition, my engagement with campus research leadership through national initiatives and grant-funded coordination activities enhanced Kansas State University’s presence in cross-institutional research infrastructure and data governance efforts.

\subsection*{Service to the Library Profession and National Communities}

At the professional and national levels, I provided sustained service through leadership, editorial work, committee participation, and grant-funded coordination activities.

I founded and served as Editor (and Co-Editor-in-Chief) of \textit{Cyberbiosecurity Quarterly}, a peer-reviewed scholarly journal hosted on New Prairie Press. In this role, I established editorial governance, developed peer-review workflows, and helped launch a new scholarly venue at the intersection of genomics, cybersecurity, and biosecurity. The journal has received public recognition from partner organizations, including BIO-ISAC.

I served as a Co-Principal Investigator on the National Science Foundation Research Coordination Network project \textit{``Reimagining a Sustainable Data Network to Accelerate Agricultural Research and Discovery''} (Award \#2126334). Through this role, I contributed to national coordination of community workshops, working groups, and annual meetings supporting data standards, interoperability, and scholarly communication in agricultural genomics.

I also served on the AgBioData Steering Committee and the AgBioData Scientific Literature Working Group, contributing to national governance and standards development for agricultural genomics data. In addition, I continued service with the i5k initiative through the Standards Working Group and assumed leadership of the i5k Training Working Group.

I joined the advisory board of the Bioeconomy Information Sharing and Analysis Center (BIO-ISAC) and participated in the BIO-ISAC annual conference, contributing to national dialogue on cyberbiosecurity, data governance, and information sharing across the bioeconomy.

\subsection*{Service to the Public}

My non-directed service also supported public access to scholarship and broader community engagement through open publishing and open infrastructure initiatives. Through the establishment of \textit{Cyberbiosecurity Quarterly} and participation in FORCE11 scholarly communications training, I contributed to advancing open, transparent, and community-governed models of scholarly communication that benefit researchers, practitioners, and the public.

\subsection*{2026 Professional Contributions}
\begin{itemize}
  \item Received the AgBioData Mission and Leadership Award in April 2026 (\href{https://experts.ksu.edu/david.molik/professional}{K-State Experts service record}); the exact award date differs between the profile and achievement log.
  \item Completed four article reviews and six meta-reviews for the First International Workshop on Cyberbiosecurity, recorded March 27. These are ten review assignments, not necessarily ten distinct manuscripts.
  \item Served as recommender for \href{https://doi.org/10.24072/pci.genomics.100695}{\textit{ViroSeek: a viral detection pipeline for second-generation sequencing}}, with the Round 2 recommendation dated August 11. This is editorial/peer-review service rather than authorship of the underlying research article.
  \item Served on the 2026 Ad Hoc Committee to Clarify Tenure and Promotion Procedures, as recorded in the supplied Libraries procedures' committee list.
\end{itemize}

\section{Non-Directed Service Evidence}
\SectionGuidance{Provide supporting references or links for your service contributions. Use evidence that illustrates your role and outcomes.}
\printbibliography[keyword={nondirected},heading=none]
